\documentclass[10pt,a4paper]{amsart}
\usepackage[margin=19mm]{geometry}
\usepackage{amsmath,amssymb,mathtools}
\usepackage[scr=rsfs,cal=euler]{mathalfa}
\usepackage{xfrac}
\usepackage[unicode,hidelinks,bookmarks=false]{hyperref}
\newcommand{\ie}{i.e.\ }
\newcommand{\cf}{cf.\ }
\newcommand{\resp}{respectively\ }
\newcommand{\Subsection}{\subsection}
\providecommand{\nobreakdash}{\nobreak\hbox{-}}
\setlength{\emergencystretch}{2em}
\multlinegap0pt
\usepackage{mathtools}
\usepackage{amssymb}
\usepackage{xcolor}
\usepackage[normalem]{ulem}
\usepackage[shortlabels]{enumitem}
\newtheorem{lemma}{Lemma}
\newcommand{\ind}{\operatorname{ind}}
\title{Finite abelian subgroups of algebraic groups}
\author{Danny Ofek, Zinovy Reichstein, Federico Scavia}
\date{Source: arXiv:2608.01595v1 (CC BY 4.0)}
\begin{document}
\maketitle

\noindent\emph{Source and licence.} Finite abelian subgroups of algebraic groups. Version arXiv:2608.01595v1 (CC BY 4.0), distributed by arXiv under the Creative Commons Attribution 4.0 International licence, http://creativecommons.org/licenses/by/4.0/, as recorded in the arXiv licence metadata for this exact version and shown on the abstract page https://arxiv.org/abs/2608.01595v1. Full licence notice: \texttt{licenses/arxiv-2608.01595-CC-BY-4.0.txt}.\par\smallskip

\noindent\textit{Editorial scope.} Counted unit: the article's lemma lem.t(G) (source0.tex lines 445--448). The article's complete original proof of it is delivered with it (source0.tex lines 450--472, one \texttt{proof} environment of 23 lines). No mathematics is written, repaired or re-derived: the statement and the proof are the article's own lines, in the article's own notation, and no retained line is substituted or re-flowed.  No further source-local context is needed: every cross-reference of every retained line resolves inside the two delivered spans.  Nothing was omitted from any retained span, so the package carries no omission record.  No retained line contains a citation, so the wrapper carries no bibliography and no bibliography entry.  No retained line contains a figure environment, an image-inclusion command or drawing code, so no image is delivered and the package declares no asset dependency.  Editorial formatting only, in addition: an AMS \texttt{amsart} preamble that restates the article's own declarations and macros used by the retained lines, namely the environments \texttt{lemma} on separate counters, together with the standard packages those lines need.  The retained environments print in this document as Lemma 1 in order of appearance, whereas the article prints the same items with the numbering of its own full document.  The complete native source of the article as distributed in the arXiv e-print for this exact version is bundled unchanged as \texttt{sources/206643/source0.tex}.
\begin{lemma} \label{lem.t(G)}
Let 
\[1 \longrightarrow G_1 \longrightarrow G \longrightarrow G_2 \longrightarrow 1\] be a short exact sequence of affine algebraic $k$-groups. Then $t(G)$ divides $t(G_1) t(G_2)$.
\end{lemma}

\begin{proof} Let $K$ be a field containing $k$, and let $E$ be a $G$-torsor over $K$.  
Letting $E_2=E/G_1$ be the $G_2$-torsor over $K$ induced by $E$, by definition of $t(G_2)$ there exist field extensions $K_1/K, \ldots, K_m/K$ of degrees $d_i\coloneqq [K_i : K]$ such that $E_2(K_i)\neq\emptyset$ for all $i=1,\dots,m$ and 
\begin{equation} \label{e.t(G_2)}
\text{$\gcd(d_1, \ldots, d_m)$ divides $t(G_2)$.}
\end{equation}
For all $i=1,\dots,m$, since $E_2(K_i)\neq\emptyset$, the $G$-torsor $E_{K_i}$ admits reduction of structure to a $G_1$-torsor $\tilde{E}_i$ over $K_i$. By definition of $t(G_1)$, for all $i = 1, \ldots, m$ there exist field extensions $K_{i\, 1}/K_i, \ldots, K_{i\,n_i}/K_i$ of degrees $d_{ij} \coloneqq [K_{ij}: K_i]$ such that each $\tilde{E}_i(K_{ij})\neq\emptyset$, and
\begin{equation} \label{e.t(G_1)}
\text{$\gcd(d_{i\, 1}, \ldots, d_{i\, n_i})$ divides $t(G_1)$.}
\end{equation} 
For all $i,j$, since $\tilde{E}_i(K_{ij})\neq\emptyset$, we have that $E(K_{ij})\neq\emptyset$, and hence $\ind(E)$ divides  
\[ [K_{ij}: K] = [K_{ij}: K_i] \cdot [K_i: K] = d_{ij} \cdot d_i.\]
Therefore, in order to finish the proof of Lemma~\ref{lem.t(G)}, it suffices to show that the ideal $\Lambda$ of $\mathbb Z$ generated by $d_i d_{ij}$ for every $i = 1, \ldots, m$ and $j = 1, \ldots, n_i$ contains $t(G_1) t(G_2)$. To prove this, let $I \coloneqq (d_1, \ldots, d_m) $ be the ideal of $\mathbb Z$ generated by $d_1, \ldots, d_m$ and $J_i \coloneqq (d_{i\, 1}, \ldots, d_{i\, n_i})$ be the ideal of $\mathbb Z$ generated by $d_{i\, 1}, \ldots, d_{i\, n_i}$. Then $(t(G_2)) \subset I$ by~\eqref{e.t(G_2)}, and $(t(G_1)) \subset J_i$ for each $i$  by~\eqref{e.t(G_1)}. Now
\begin{align*}
(t(G_1)t(G_2))
&= (t(G_1)) \cdot (t(G_2)) \\
&\subset (t(G_1)) \cdot I \\
&= (t(G_1)) \cdot (d_1) + (t(G_1)) \cdot (d_2) + \ldots
   + (t(G_1)) \cdot (d_m) \\
&\subset J_1 \cdot (d_1) + J_2 \cdot (d_2) + \ldots + J_m \cdot (d_m) \\
&= \Lambda,
\end{align*}
as desired.     
\end{proof}

\end{document}
